\documentclass[a4paper,11pt]{article}

% -------------------------------------------------
% Sprache und Kodierung
% -------------------------------------------------
\usepackage[utf8]{inputenc}
\usepackage[T1]{fontenc}
\usepackage[ngerman]{babel}

% -------------------------------------------------
% Layout und Typografie
% -------------------------------------------------
\usepackage[a4paper,margin=2.5cm]{geometry}
\usepackage{lmodern}
\usepackage{microtype}
\usepackage{setspace}
\usepackage{parskip}

% -------------------------------------------------
% Listen und Überschriften
% -------------------------------------------------
\usepackage{enumitem}
\usepackage{titlesec}

% -------------------------------------------------
% Farben und Boxen
% -------------------------------------------------
\usepackage{xcolor}
\usepackage[most]{tcolorbox}

% -------------------------------------------------
% Kopf- und Fußzeile
% -------------------------------------------------
\usepackage{fancyhdr}
\usepackage{lastpage}

% -------------------------------------------------
% Hyperlinks
% -------------------------------------------------
\usepackage[hidelinks]{hyperref}

% -------------------------------------------------
% Einstellungen
% -------------------------------------------------

\setstretch{1.12}

\definecolor{mainblue}{RGB}{35, 82, 124}
\definecolor{lightblue}{RGB}{240, 246, 252}
\definecolor{darkgray}{RGB}{70, 70, 70}

\titleformat{\section}
{\large\bfseries\color{mainblue}}
{}
{0pt}
{}

\titlespacing*{\section}
{0pt}
{1.2em}
{0.5em}

\setlist[itemize]{
	leftmargin=1.2cm,
	itemsep=0.25em,
	topsep=0.4em
}

\pagestyle{fancy}
\fancyhf{}
\lhead{\small Beschreibung der Ist-Situation}
\rhead{\small Lieferbestätigung}
\cfoot{\small Seite \thepage\ von \pageref{LastPage}}

\renewcommand{\headrulewidth}{0.4pt}
\renewcommand{\footrulewidth}{0pt}

\newtcolorbox{infobox}{
	colback=lightblue,
	colframe=mainblue,
	arc=2mm,
	boxrule=0.5pt,
	left=8pt,
	right=8pt,
	top=7pt,
	bottom=7pt
}

% -------------------------------------------------
% Dokument
% -------------------------------------------------

\begin{document}
	
	% -------------------------------------------------
	% Titelbereich
	% -------------------------------------------------
	
	\begin{center}
		\vspace*{0.5cm}
		
		{\LARGE\bfseries\color{mainblue}
			Beschreibung der aktuellen Ist-Situation}
		
		\vspace{0.3cm}
		
		{\large Lieferbestätigung im aktuellen Dispositionsprozess}
		
		\vspace{0.5cm}
		
		\rule{0.75\textwidth}{0.5pt}
		
		\vspace{0.3cm}
		
		{\small Stand: \today}
	\end{center}
	
	\vspace{0.8cm}
	
	\begin{infobox}
		Dieses Dokument beschreibt den aktuellen manuellen Prozess zur telefonischen
		Lieferbestätigung im Rahmen der Disposition. Der Fokus liegt auf der bestehenden
		Ist-Situation, den beteiligten Rollen, den Entscheidungsfällen sowie der
		systemseitigen Dokumentation.
	\end{infobox}
	
	\vspace{0.8cm}
	
	% -------------------------------------------------
	% Inhalt
	% -------------------------------------------------
	
	\section*{1. Ausgangssituation}
	
	Der eigentliche Prozess beginnt erst, nachdem ein Auftrag bereits im System
	vorhanden ist. Die Bestellung erfolgt zuvor über bestehende Kanäle,
	beispielsweise über Online-Plattformen wie HeizOel24.
	
	Für die interne Bearbeitung liegen anschließend bereits folgende Informationen vor:
	
	\begin{itemize}
		\item Kunde
		\item Auftrag
		\item Lieferadresse
		\item Produkt
		\item Bestellmenge
	\end{itemize}
	
	Zusätzlich wurde zuvor vom Kunden bereits ein Lieferzeitraum angegeben,
	zum Beispiel ein Zeitraum wie 01.04.--06.04. mit Einschränkungen wie
	\glqq nicht nach 16 Uhr\grqq.
	
	\section*{2. Rolle der Disposition}
	
	Der Disponent erstellt auf dieser Grundlage die vorläufige Tourenplanung.
	Eine Tour entspricht dabei dem Lieferweg eines einzelnen Fahrzeugs.
	
	Dabei plant der Disponent die Tour im Detail, erstellt eine Route und schlägt
	innerhalb dieser Route optimale Lieferzeitpunkte vor.
	
	Manuell festgelegt werden insbesondere:
	
	\begin{itemize}
		\item Liefertag
		\item konkretes Lieferfenster
	\end{itemize}
	
	Das Lieferfenster wird zunächst intern disponiert. Der vorgeschlagene
	Lieferzeitpunkt muss anschließend telefonisch mit dem Kunden abgestimmt werden.
	
	\section*{3. Telefonische Rückbestätigung}
	
	Spätestens zwei Tage vor dem geplanten Liefertermin erfolgt ein telefonischer
	Kontakt durch den Disponenten.
	
	Dem Kunden wird der Lieferauftrag mit den vereinbarten Informationen genannt,
	insbesondere:
	
	\begin{itemize}
		\item Preis pro 100 Liter
		\item Lieferadresse
		\item Lieferdatum
		\item geplantes Lieferfenster
	\end{itemize}
	
	Ziel ist die Klärung, ob der konkrete Liefertermin bestätigt werden kann.
	
	Ein Kontakt gilt nur dann als erfolgreich, wenn eine konkrete Uhrzeit
	beziehungsweise ein eindeutiges Lieferfenster bestätigt wird. Aussagen wie
	eine spätere Rückmeldung oder eine unklare Zusage gelten nicht als Bestätigung.
	
	\section*{4. Ablauf bei Nichterreichbarkeit}
	
	Wenn der Kunde beim ersten Anruf nicht erreichbar ist, erfolgt am selben Tag
	nach einigen Stunden ein weiterer Anruf.
	
	Je nach Situation können zusätzlich ein dritter oder vierter Anruf erfolgen.
	
	Falls weiterhin kein Kontakt zustande kommt:
	
	\begin{itemize}
		\item bleibt der Auftrag grundsätzlich weiterhin in der Planung,
		\item da der Kunde bereits einen Lieferzeitraum angegeben hat,
		\item und die Tour auf dieser Basis weiter disponiert wird.
	\end{itemize}
	
	Falls der Kunde später selbst zurückruft, prüft der Disponent, ob innerhalb
	der ursprünglichen Tour noch freie Kapazität vorhanden ist.
	
	Ist das Zeitfenster bereits vergeben, erfolgt eine Verschiebung auf die nächste
	freie Tour.
	
	\section*{5. Ablauf bei Kundenreaktion}
	
	Wenn der Kunde den vorgeschlagenen Termin bestätigt, gilt dieser als verbindlich.
	
	Der Auftrag bleibt final in der geplanten Tour. In diesem Zustand gilt der
	Auftrag als fest eingeplant und ist technisch für weitere Änderungen gesperrt.
	
	Dem Kunden kann in diesem Zusammenhang eine Anfahrtspauschale berechnet werden.
	Falls der Kunde zum bestätigten Termin nicht anwesend ist, entsteht eine
	Vertragsstrafe.
	
	Falls der Kunde einen bereits bestätigten Termin kurzfristig absagt, wird der
	Auftrag aus der aktuellen Tour entfernt und auf eine spätere freie Tour
	verschoben.
	
	Spontane Absagen nach Bestätigung können zusätzlich zu verrechneten
	Anfahrtspauschalen oder Planungspauschalen führen.
	
	\section*{6. Änderungswünsche des Kunden}
	
	Falls der Kunde den vorgeschlagenen Termin nicht akzeptiert, wird geprüft,
	ob innerhalb der aktuellen Tour ein alternatives Zeitfenster verfügbar ist.
	
	Typische Fälle sind:
	
	\begin{itemize}
		\item Lieferung früher gewünscht
		\item Lieferung später gewünscht
		\item Einschränkungen wie nur vormittags oder erst nach 16 Uhr
	\end{itemize}
	
	Teilweise wurden diese Einschränkungen bereits bei der Auftragserteilung vom
	Kunden angegeben und an den Disponenten übergeben.
	
	Ist das gewünschte Zeitfenster innerhalb der aktuellen Tour verfügbar, wird der
	Auftrag entsprechend angepasst.
	
	Führt die Änderung zu Konflikten im logistischen Ablauf der Tour, entscheidet
	der Disponent im Einzelfall über Aufnahme oder Verschiebung des Auftrags.
	
	Eine formalisierte Entscheidungsregel existiert hierfür nicht.
	
	Ist keine Anpassung möglich, erfolgt eine Verschiebung auf eine andere freie Tour.
	
	\section*{7. Alternative Kontaktwege}
	
	In Ausnahmefällen erfolgt die Kontaktaufnahme manuell über alternative
	Kommunikationskanäle.
	
	Verwendet werden dabei insbesondere:
	
	\begin{itemize}
		\item Handynummer
		\item Telefonnummer
		\item E-Mail
		\item WhatsApp
	\end{itemize}
	
	Auch hierbei erfolgt die Bearbeitung vollständig durch den Disponenten.
	
	\section*{8. Systemseitige Dokumentation}
	
	Alle Rückmeldungen werden im DISPO-System dokumentiert.
	
	Dort existieren Statusinformationen zum Bearbeitungsstand des Auftrags.
	
	Nichterreichbarkeit ergibt sich dabei aus dem fehlenden bestätigten Rückkontakt.
	
	\section*{9. Hauptcharakteristik des aktuellen Prozesses}
	
	Der aktuelle Prozess basiert vollständig auf manueller Abstimmung durch die
	Disposition.
	
	Zentrale Entscheidungen bei Terminbestätigung, Verschiebung und Sonderfällen
	werden individuell durch den Disponenten getroffen.
	
	Ein automatisierter digitaler Workflow für Rückbestätigung, Statusverarbeitung
	oder Folgeaktionen existiert derzeit nicht.
	
	Der Prozess endet, sobald der Liefertermin bestätigt, der Auftrag verschoben oder keine weitere Rückmeldung des Kunden vor dem geplanten Liefertermin möglich ist.
	\vfill
	
	\begin{center}
		\rule{0.6\textwidth}{0.4pt}
		
		\vspace{0.2cm}
		
		{\small Ende der Beschreibung der aktuellen Ist-Situation}
	\end{center}
	
\end{document}